\section{Removing a fixed pseudoeffective error}
\label{conv:section}

The controlled-boundary theorem becomes useful for anticanonical
nonvanishing when the boundary on an auxiliary base disappears after
pullback and birational pushforward.  We now construct precisely such a
base.  The error in the following theorem is an arbitrary fixed
pseudoeffective divisor; it need not be effective.

\begin{theorem}\label{conv:fixed-error}
Let $X$ be a smooth connected projective complex variety and let
$L=-K_X$ have a smooth semipositive Hermitian metric.  Let $P$ be a
pseudoeffective Cartier divisor.  If there are effective integral
$N_j\sim m_jL-P$ for unbounded positive integers $m_j$, then
$H^0(X,mL)\ne0$ for some positive integer $m$.
\end{theorem}

We use the Iitaka fibration of just one $N_j$.  Ratios of its sections
force all subsequent divisors to vary affinely on the generic fiber.
The vertical slopes define a pseudoeffective divisor on the base.  A
rank-one adjoint direct image supplies the metric required by
Theorem~\ref{thm:controlled-boundary}, and its boundary is exceptional
over $X$.  Thus the correction allowed by that theorem causes no poles
in the section returned to $X$.

\subsection{Vertical minima and return of sections}

We record the return step once, including the exceptional divisors that
can appear when a smooth model ceases to be equidimensional.

\begin{lemma}\label{conv:return}
Let $\pi:Z\to X$ be a projective birational morphism of smooth
projective varieties and let $f:Z\to Y$ be a surjective morphism to a
smooth projective variety with connected general fibers, and let $L$ be
a rational line bundle on $X$.  Assume every
prime of $Z$ mapping into codimension at least two in $Y$ is
$\pi$-exceptional.  Let $R\sim_{\mathbb Q}a\pi^*L$, $a>0$, be a
rational divisor with nonnegative horizontal part.  For each base prime
$Q$ put
\[
 \mu_Q(R)=\min_{D\mapsto Q}
 \frac{\operatorname{ord}_D R}{\operatorname{ord}_D f^*Q},
 \qquad M=\sum_Q\mu_Q(R)Q.
\]
Suppose $M$ is pseudoeffective.  If there is an effective rational
$C_0$, all components of whose pullback are $\pi$-exceptional, such that
$M+C_0\sim_{\mathbb Q}G\ge0$, then $L$ is rationally linearly
equivalent to an effective rational divisor.  The same conclusion holds
if the premise is $\epsilon M+C_0\sim_{\mathbb Q}G\ge0$ for some
rational $\epsilon>0$.

For the invariant conclusion, give $L$ a specified linearization and
assume
\[
 R=c\,\operatorname{div}(\pi^*t),\qquad cr=a,
\]
where $c>0$ is rational and $t$ is a nonzero invariant rational section
of the specified linearized $L^r$, with $r>0$ an integer for which
$L^r$ is a line bundle.  Assume also that every function pulled back
from $Y$ is invariant in $\mathbb C(X)$.  Then the resulting section
of a positive multiple of $L$ is invariant for that linearization;
no regular group action on $Z$ is required.
\end{lemma}
\begin{proof}
Only finitely many coefficients of $M$ are nonzero.  Replace $C_0,G$ by
$C_0/\epsilon,G/\epsilon$ in the last variant.  The divisor
\[
 R+f^*(G-M-C_0)\sim_{\mathbb Q}a\pi^*L
\]
has nonnegative horizontal coefficients.  At a prime dominating $Q$
outside the support of $C_0$, this follows from the minimum defining
$M$ and from $G\ge0$.  All other possibly negative primes are
$\pi$-exceptional, either by the support premise on $C_0$ or by the
assumption on primes mapping into codimension two.  Its pushforward is
therefore effective and rationally linearly equivalent to $aL$.
Clearing denominators proves the assertion.  In the invariant case
choose the clearing integer $n$ also so that $nc$ is integral.  The
operation multiplies $t^{nc}$, a section of $L^{ncr}=L^{na}$, by a
rational function pulled from $Y$.  Both factors are invariant, and
extension across codimension two on the smooth variety $X$ preserves
invariance.
\end{proof}

The required model condition follows from flattening, rather than from
an assertion that a final smooth resolution is flat.  Starting with a
resolved rational map and its Stein factorization, flatten over a
modification of the base \cite[Theorem~5.2.2]{RaynaudGruson1971}.
Normalize the main component, resolve the base, flatten again if
necessary, and then resolve the source.  On the intervening
flat model a prime cannot map into base codimension two: its inverse
image has dimension at most $\dim X-2$.  Any prime on the final
resolution with that property is consequently exceptional over the
intervening model, and hence over $X$.

\subsection{The Iitaka base and its strict metric}
\label{conv:iitaka}

\begin{proof}[Proof of Theorem~\ref{conv:fixed-error}]
Pass to a subsequence with $m_0<m_1<m_2<\cdots$, and take the Iitaka
fibration of $N_1$.  Resolve a system of maximal image dimension, take
its Stein factorization, and prepare the model by flattening as above:
\[
 X\xleftarrow{\ \pi\ }Z\xrightarrow{\ f\ }Y.
\]
The base field $\mathbb C(Y)$ is relatively algebraically closed in
$\mathbb C(X)$ and contains every ratio of sections of every positive
multiple of $N_1$.  Indeed, a ratio transcendental over the original
image field would increase the image dimension after multiplication of
sections into a common system.  All other ratios lie in its relative
algebraic closure, which is the Stein field.  If $\dim Y>0$, we may
choose an ample $H$ on $Y$ and an integer $b>0$ with
\begin{equation}\label{conv:iitaka-polarization}
 b\pi^*N_1\sim f^*H+D_1,\qquad D_1\ge0.
\end{equation}
The polarization of the original image pulls back to a big divisor on
$Y$; a sufficiently large multiple dominates $H$, which proves this
assertion on the modified base.

Put $R=(N_1-N_0)/(m_1-m_0)\sim_{\mathbb Q}L$.  For $j>1$ the two
effective divisors
\[
 (m_j-m_1)N_0+(m_1-m_0)N_j,
 \qquad (m_j-m_0)N_1
\]
are linearly equivalent.  Their ratio is a base function.  Consequently
there are rational principal divisors $Q_j\sim_{\mathbb Q}0$ on $Y$
with the exact identities
\begin{equation}\label{conv:affine-divisors}
 \pi^*N_j=\pi^*N_0+(m_j-m_0)(\pi^*R+f^*Q_j).
\end{equation}
Effectivity as $m_j\to\infty$ makes the horizontal part of $\pi^*R$
nonnegative.  If $Y$ is a point, this already proves the theorem.  If
$\dim Y=\dim X$, then $N_1$ is big, so $m_1L\sim N_1+P$ is big and
again has a section in a positive multiple.  We henceforth assume
$0<\dim Y<\dim X$.

Define $M=\sum_Q\mu_Q(\pi^*R)Q$ as in
Lemma~\ref{conv:return}.  There is a fixed effective rational $J$ on
$Y$ such that
\[
 M+Q_j+\frac{J}{m_j-m_0}\ge0.
\]
For example its coefficient at $Q$ can be the maximum of
$\operatorname{ord}_D\pi^*N_0/\operatorname{ord}_D f^*Q$ over
$D\mapsto Q$.  These maxima have finite support.  The inequality
follows by taking each coefficient in \eqref{conv:affine-divisors}.
Since $Q_j$ is rational principal, the numerical classes of these
effective divisors converge to $[M]$.  Thus $M$ is pseudoeffective.

The remaining problem is to make $M$ effective with a correction
exceptional over $X$.  Write $E=K_{Z/X}\ge0$.  We first show
\begin{equation}\label{conv:adjoint-rank}
 \operatorname{rank}f_*\mathcal O_Z(E)=1.
\end{equation}
The Jacobian section supplies rank at least one.  If the rank were
larger, twisting by $kH$ for large $k$ would supply two sections of
$E+kf^*H$ independent on the generic fiber.  Multiply by the section of
$bk\pi^*N_1-kf^*H$ furnished by
\eqref{conv:iitaka-polarization}.  The resulting ratio belongs to
$|E+bk\pi^*N_1|$.  Since
$\pi_*\mathcal O_Z(E)=\mathcal O_X$, it is a ratio from a multiple of
$N_1$, and hence a base function, a contradiction.  The equality of
pushforwards used here follows from normality of $X$: a rational
function with possible poles only on exceptional divisors has no pole
at any prime of $X$ and extends across codimension two.

Write
\[
 (f_*\mathcal O_Z(E))^{**}=\mathcal O_Y(C),\qquad C\ge0,
\]
where $C$ is the divisor of the canonical section.  The divisor $C$ is
integral.  The torsion-free rank-one direct image agrees with its hull
at every codimension-one point.  If $Q\subset\operatorname{Supp}C$,
dividing the canonical section by a local equation of $Q$ is therefore
allowed at its generic point.  Hence every prime of $Z$ dominating
$Q$ occurs in $E$.  The model condition handles components over higher
codimension, so every component of $f^*C$ is exceptional over $X$.

We now verify all analytic premises of
Theorem~\ref{thm:controlled-boundary}.  Combining
\eqref{conv:iitaka-polarization} with $N_1\sim m_1L-P$ gives
\[
 bm_1\pi^*L\sim b\pi^*P+f^*H+D_1.
\]
A positive-current metric for $P$, an ample metric for $H$, and the
divisor metric of $D_1$ therefore give a singular metric on $\pi^*L$
whose curvature dominates a positive multiple of $f^*\omega_Y$.
Mix a sufficiently small positive fraction of its weight with the
smooth semipositive pullback metric.  Skoda's small-exponent
integrability theorem, on a finite cover of the compact $Z$, gives a
metric $h$ satisfying
\begin{equation}\label{conv:strict-upstairs}
 \Theta_h\ge\delta f^*\omega_Y,
 \qquad \mathcal J(h)=\mathcal O_Z,
 \qquad \delta>0.
\end{equation}
The metric $h$ need not be smooth.  The triviality of its multiplier
ideal is a statement on the whole total space; a uniform claim on
all fibers is neither needed nor asserted.

Apply singular adjoint direct-image positivity
\cite[Theorem~3.3.5]{PT2018}; see also
\cite[Theorem~21.1 and Corollary~21.2]{HPS2018}.
The natural inclusion of the multiplier-ideal direct image into the
ordinary adjoint direct image is an equality here, because
$\mathcal J(h)=\mathcal O_Z$ in \eqref{conv:strict-upstairs}.
In particular the generic-isomorphism hypothesis in the cited theorem
holds, and it gives a metric on the ordinary sheaf.  Projection formula and
\eqref{conv:adjoint-rank} identify its rank-one hull as
\[
 \bigl(f_*\mathcal O_Z(K_{Z/Y}+\pi^*L)\bigr)^{**}
   =\mathcal O_Y(-K_Y+C).
\]
It acquires a metric $g$ with $\Theta_g\ge\delta\omega_Y$.
For the lower bound, subtract a local potential of
$\delta f^*\omega_Y$ upstairs; the fiber integral subtracts the
corresponding base potential downstairs.  The direct-image theorem
extends the metric across divisors; plurisubharmonic extension handles
the remaining codimension-two set.  Thus this is a global strict
curvature bound on the actual line bundle, including its divisorial
correction.

The distinguished section has finite norm.  On the smooth base-change
locus its norm times coordinate base volume is the pushforward of the
squared norm of the Jacobian section of $K_Z+\pi^*L$.
The latter is locally integrable by $\mathcal J(h)=\mathcal O_Z$.
Properness and Fubini give finite local mass on the base; the omitted
analytic subsets have measure zero.  In the line $-K_Y+C$ the
distinguished local section is $s_C$.  Consequently
\[
 \mathcal J(g)\supset\mathcal O_Y(-C).
\]
We have verified global strictness, the precise multiplier-ideal
inclusion, and exceptional pullback of the integral boundary.
Theorem~\ref{thm:controlled-boundary} gives
$M+C_0\sim_{\mathbb Q}G\ge0$, with $C_0\ge0$ supported on $C$.
Lemma~\ref{conv:return} now gives the required section on $X$.
\end{proof}

\subsection{The smooth toroidal metric for the same base}
\label{conv:toroidal-alternative}

A toroidal embedding is locally modeled on a toric variety with its
dense torus; its boundary corresponds to the toric boundary.  We use
strict toroidal embeddings, whose boundary components are normal.
A toroidal morphism is locally a toric morphism, described on the dense
tori by monomials.  This local description keeps track of the exact
ramification exponents in the metric correction.

There is another way to supply the metric on the Iitaka base.  It keeps
the original smooth anticanonical metric throughout fiber integration
and gives unweighted integrability.  We give the geometric verification
because the two metric constructions have different conclusions.

\begin{lemma}\label{conv:null-curvature}
Let $Z$ be smooth projective of dimension $n$, and let $\theta,\eta$ be
smooth semipositive closed $(1,1)$-forms.  Suppose $a[\theta]-[\eta]$
is pseudoeffective for some $a>0$.  If $\nu$ is the maximum rank of
$\theta$, then $\ker\theta\subset\ker\eta$ on its rank-$\nu$ locus.
If $\eta=f^*\omega_Y$ on an ordinary open meeting the rank-$\nu$
locus, with $f$ submersive and $\omega_Y>0$, there is a smaller
nonempty patch on which
$\theta\ge c f^*\omega_Y$ for some $c>0$.
\end{lemma}
\begin{proof}
There is nothing to prove when $\nu=n$.  Otherwise choose a K\"ahler
form $\omega$.  Every mixed intersection of $\theta$ and $\eta$ of
total degree $\nu+1$, completed by $\omega^{n-\nu-1}$, is
nonnegative.  Replacing one $\eta$ at a time by $a\theta$ can only
increase the intersection, because a pseudoeffective class intersects
nonnegatively with $n-1$ nef classes.  The resulting upper bound is a
multiple of $\int_Z\theta^{\nu+1}\omega^{n-\nu-1}=0$.
Thus $(\theta+\eta)^{\nu+1}=0$ pointwise, by semipositivity.  On the
rank-$\nu$ locus $\theta+\eta$ has the same kernel as $\theta$;
this is the asserted inclusion.  On a relatively compact constant-rank
patch, the nonzero eigenvalues of $\theta$ have a positive lower bound
and $f^*\omega_Y$ is bounded, giving the final inequality.
\end{proof}

We also isolate the algebraic consequence of unweighted integrability.

\begin{corollary}\label{conv:unweighted-effectivity}
Let $Y$ be smooth projective of positive dimension and $A_0\ge0$
rational.  Suppose $B=-K_Y+A_0$ has a semipositive metric with locally
integrable squared frame norms, smooth and strictly positive on a
nonempty ordinary open.  For every pseudoeffective rational $M$ there
is a rational $\epsilon>0$ such that
$A_0+\epsilon M$ is rationally linearly equivalent to an effective
rational divisor.
\end{corollary}
\begin{proof}
Apply Lemma~\ref{lem:integrable-neighborhood} in the direction $M$.
It gives an ample rational $H'$, a sufficiently small rational
$\epsilon>0$, and an effective klt boundary
$\Delta_0\sim_{\mathbb Q}B+\epsilon M-H'$.
Represent $H'$ by a general effective rational divisor with small
coefficients transverse to a log resolution of $\Delta_0$.
Its addition gives an effective big klt boundary
$\Delta\sim_{\mathbb Q}B+\epsilon M$.  The log canonical divisor
$K_Y+\Delta\sim_{\mathbb Q}A_0+\epsilon M$ is pseudoeffective.
Big-boundary nonvanishing \cite[Theorem~D]{BCHM2010} applies.
Lemma~\ref{lem:rational-recovery} gives rational linear effectivity.
\end{proof}

Prepare the rational Iitaka map birationally as an equidimensional
strict toroidal morphism $f:W\to Y$, with $W$ normal, $Y$ smooth,
and $p:W\to X$ birational.  Such a preparation follows from weak
toroidalization \cite[Theorem~1.1]{ADK2013} and equidimensional
toroidal subdivision \cite[Proposition~4.4]{AK2000}.  Mark the
nonisomorphism locus in the source boundary before these operations.
No finite base alteration or reduced-fiber assertion is used here.
Take a smooth resolution $Z\to W$, with maps $\pi:Z\to X$ and
$g:Z\to Y$.  The resolution itself need not be equidimensional.
The generic rank-one argument \eqref{conv:adjoint-rank} is unchanged.
Likewise the horizontal interpolation and the vertical minima may be
computed on $W$, where every vertical prime dominates a base prime.

Equation~\eqref{conv:iitaka-polarization}, pulled to $Z$, says that a
multiple of $[\Theta_{\pi^*L}]$ dominates the pullback of a positive
base class pseudoeffectively.  Lemma~\ref{conv:null-curvature} gives a
horizontal strictness patch above the smooth base-change locus, away
from the Jacobian divisor.  The Jacobian section is nonzero there.
The model, rank-one, and strict-patch hypotheses have now been
verified for the original smooth coefficient metric.  These are the
hypotheses of the toroidal-volume theorem in
\cite[Theorem~2.1]{CompanionG}.
That theorem supplies a rational $A_0\ge0$ with
$p_*f^*A_0=0$ and the unweighted integrable metric of
Corollary~\ref{conv:unweighted-effectivity}.  Its coefficients are
\[
 \operatorname{ord}_Q A_0=
 \max\left\{0,\min_{D\mapsto Q}
       \left(\frac{1+\operatorname{ord}_D K_{W/X}}
                        {\operatorname{ord}_D f^*Q}-1\right)\right\}.
\]
The all-valuation integrability and extension assertions here are part
of that theorem, rather than consequences of generic rank one alone.
Corollary~\ref{conv:unweighted-effectivity} and
Lemma~\ref{conv:return}, on a common resolution, complete the same
conversion.  The output here is unweighted integrability for a rational
correction: a nonvanishing frame has finite norm.  The smooth coarea
construction in Section~\ref{coa:analytic-tools} instead gives an
integral correction and an $L^{1+\eta}$ bound for the distinguished
section.  Lemma~\ref{coa:weighted-effectivity} converts that output
to the same supported effectivity premise in
Lemma~\ref{conv:return}.

\subsection{The field of all pseudoeffectively dominated systems}
\label{conv:all-dominated}

A second base construction uses every effective system dominated by a
multiple of $L$.  It is useful when no distinguished divisor should be
chosen in advance.  We retain the different maximality argument and the
explicit integral correction it produces.

Call a subsystem of an integral divisor $D$ admissible if it has a
nonzero section and $bL-D$ is pseudoeffective for some integer $b>0$.
Choose finitely many admissible systems whose joint image has maximum
dimension $d$.  Every further admissible ratio is algebraic over this
image field.  Take its relative algebraic closure and a smooth
projective connected-fiber model $\pi:Z\to X$, $f:Z\to Y$, prepared
by flattening as before.  For every ample $H$ on $Y$ there is a $b_0>0$
with
\begin{equation}\label{conv:all-polarization}
 b_0\pi^*L-f^*H\quad\hbox{pseudoeffective}.
\end{equation}
Indeed the sum of the chosen divisors dominates the pullback of the
product polarization; its restriction to the joint image is big.
The sum itself is pseudoeffectively dominated by a multiple of $L$,
and a large multiple of the product polarization dominates $H$.
If $d=\dim X$, this makes $L$ big.

Take sections $u_j$ defining $N_j\sim m_jL-P$, put
$r=m_2-m_1$, and let $t=\pi^*(u_2/u_1)$, a rational section of
$r\pi^*L$.  For $j>2$ the expressions
$u_j^r u_1^{m_j-m_2}$ and $u_2^{m_j-m_1}$ belong to the same
admissible system, namely a multiple of $m_2L-P$.
Writing $D_t=\operatorname{div}(t)$ gives
\[
 r\pi^*N_j=r\pi^*N_1+(m_j-m_1)D_t+
                    f^*\operatorname{div}(g_j),
 \qquad g_j\in\mathbb C(Y)^*.
\]
Thus $D_t$ has nonnegative horizontal part, and its divisor of vertical
minima $M$ is pseudoeffective by the same fixed-error inequality used
in \eqref{conv:affine-divisors}.  The case $d=0$ is already settled.

For $d>0$, the full adjoint sheaf $f_*\mathcal O_Z(E)$,
$E=K_{Z/X}$, has rank one.  Otherwise two generic-fiber independent
sections can be made global in $E+kf^*H$.  After pushforward their
effective divisors give a subsystem of $k\pi_*f^*H$ on $X$ whose
ratio is not in $\mathbb C(Y)$.  By
\eqref{conv:all-polarization} this system is admissible, contradicting
maximality.  Mixing a positive-current metric furnished by
\eqref{conv:all-polarization} with the smooth pullback metric again
gives a globally defined coefficient $h$ with
$\Theta_h\ge\delta f^*\omega_Y$ and
$\mathcal J(h)=\mathcal O_Z$.  The singular direct-image theorem
therefore applies to the full actual sheaf as in the first proof.
Let $w$ be the pushed-forward Jacobian volume density on its smooth
base-change locus.  Then $-\log w$ has curvature at least
$\delta\omega_Y$ there, and $w$ has finite local mass.

Here the boundary can be written directly.  Set
\[
 c_Q=\min_{D\mapsto Q}
       \left\lfloor\frac{\operatorname{ord}_D E}
                            {\operatorname{ord}_D f^*Q}\right\rfloor,
 \qquad C=\sum_Qc_QQ.
\]
It is effective and integral.  If a prime over $Q$ is nonexceptional
then its coefficient in $E$ is zero, so $c_Q=0$.  Hence $f^*C$ is
exceptional, including the primes over codimension two by the model
condition.  Near a general point of a prime $D$ attaining the minimum,
write $f=(z_1^e,z_2,\ldots,z_d)$ and
$a=\operatorname{ord}_D E$.  A plurisubharmonic coefficient weight is
bounded above on a relatively compact chart.  Its metric norm is
therefore bounded below by a positive constant times a smooth norm.
Integration of this chart gives
\[
 w(y)\ge c|y_1|^{2(a+1-e)/e}.
\]
If $a=qe+s$, $0\le s<e$, then
$(a+1-e)/e=q+(s+1-e)/e\le q=\lfloor a/e\rfloor$.
It follows that
$\psi=-\log w+\log|s_C|^2$ is locally bounded above at general
points of every missing divisor.  Extend plurisubharmonically across
these divisors and then across codimension two, subtracting local
potentials of $\delta\omega_Y$ to preserve the lower bound.
This gives a metric on $-K_Y+C$ with global strict curvature and
\[
 |s_C|^2e^{-\psi}=w\in L^1_{\mathrm{loc}}.
\]
Theorem~\ref{thm:controlled-boundary} applies to $M$.
Lemma~\ref{conv:return}, now with the rational section $t$, completes
the alternative construction.  Thus its distinct ingredients are the
field of all admissible systems, full rank one from maximality, and the
explicit floor correction; the algebraic return is unchanged.

\subsection{From twisted tensors to a fixed error}
\label{conv:tensors}

The determinant method of Lazi\'c--Peternell
\cite[Lemma~4.1]{LP2018}, and its formulation in
\cite[Lemma~5.1]{LMPTX2023}, provides the link to twisted differentials.
Those cited extraction statements include a numerical-nontriviality
hypothesis.  The following direct construction also covers a numerically
trivial $L$, because the starting exponents here are positive and
unbounded.

\begin{lemma}\label{conv:determinant}
Let $X$ be smooth connected projective, let $L=-K_X$ be nef, and fix
$p>0$.  If
$H^0(X,(\Omega_X^1)^{\otimes p}\otimes mL)\ne0$ for unbounded
positive integers $m$, there are a pseudoeffective Cartier divisor $P$
and effective integral $N_j\sim m_jL-P$ for unbounded positive $m_j$.
If the starting sections are invariant for a torus and the induced
linearizations, the determinant sections defining $N_j$ are invariant.
\end{lemma}
\begin{proof}
Each section defines a map $\mathcal O_X(-mL)\to
(\Omega_X^1)^{\otimes p}$.  Saturate their generic span, obtaining a
coherent subsheaf $\mathcal S$.  Finitely many maps already span at
the generic point, and saturation ensures every original map factors
through $\mathcal S$.  Let $a=\operatorname{rank}\mathcal S$.
The reflexive determinant $A=(\bigwedge^a\mathcal S)^{**}$ is a line
bundle on smooth $X$.
The cotangent-subsheaf consequence of generic tangent nefness
\cite[Theorem~4.1]{LMPTX2023} gives $-A$ pseudoeffective.
More explicitly, saturate $A$ in the corresponding exterior power of
the tensor bundle and embed that exterior power in a cotangent tensor
power.  The negative of the saturated line is pseudoeffective; passing
back to $-A$ adds an effective divisor.  Put $P=-A$.

Choose a generic basis from finitely many starting sections.  Every
nonzero section of arbitrarily large twist can replace one vector of
that basis while keeping a basis.  On an unbounded subsequence the same
slot works.  Wedge it with the remaining fixed vectors.  This gives
nonzero global sections of $A\otimes m_jL$, where $m_j$ is the
varying exponent plus a fixed sum of exponents.  The sums remain
positive and unbounded.  The wedges extend globally by the determinant
map and reflexive extension across codimension two.
If the maps are equivariant, their saturation and determinant carry
the induced action, and the wedge of invariant sections is invariant.
\end{proof}

\begin{corollary}\label{conv:tensor-conversion}
Under the metric hypotheses of Theorem~\ref{conv:fixed-error}, if
$H^0(X,(\Omega_X^1)^{\otimes p}\otimes mL)\ne0$ for a fixed
$p\ge0$ and unbounded positive $m$, then a positive multiple of $L$
has a section.  The same holds with $\Omega_X^p$ in place of the
tensor power.
\end{corollary}
\begin{proof}
For $p=0$ a section is already given.  Otherwise apply
Lemma~\ref{conv:determinant} and Theorem~\ref{conv:fixed-error}.
Exterior forms embed in tensor powers by alternation in characteristic
zero.  This conversion is an ordinary statement; invariant determinant
sections alone do not turn its conclusion into an invariant section.
\end{proof}

\subsection{A movable-slope proof of the determinant input}
\label{conv:determinant-slope}

For completeness we give a second proof of the pseudoeffectivity used
in Lemma~\ref{conv:determinant}, under its smooth-metric application
hypotheses.  This proof explains the ramification correction in the
relative-volume argument.  It uses movable slope theory
\cite{GrebKebekusPeternell2016}, duality between pseudoeffective divisors
and movable curves \cite{BDPP2013}, and the positive-minimal-slope
algebraicity theorem of Campana--P\u aun
\cite[Theorem~1.1]{CampanaPaun2019}.

For a torsion-free sheaf $\mathcal E$ and movable curve class $\alpha$,
its slope is $\mu_\alpha(\mathcal E)=c_1(\mathcal E)\cdot\alpha/
\operatorname{rank}\mathcal E$.  The symbols $\mu_{\max,\alpha}$ and
$\mu_{\min,\alpha}$ denote the greatest and least slopes in its
Harder--Narasimhan filtration, whose successive quotients are
semistable with decreasing slopes.

Suppose a saturated subsheaf of a cotangent tensor power on smooth
projective $X$ had determinant of positive degree against a movable
curve class $\alpha$.  Tensor-product slope inequalities imply
$\mu_{\max,\alpha}(\Omega_X^1)>0$, equivalently
$\mu_{\min,\alpha}(T_X)<0$.  Since $c_1(T_X)\cdot\alpha\ge0$,
the positive-slope part $\mathcal F$ of the Harder--Narasimhan
filtration of $T_X$ satisfies
\[
 \mu_{\min,\alpha}(\mathcal F)>0,
 \qquad\mu_{\max,\alpha}(T_X/\mathcal F)\le0,
 \qquad c_1(T_X/\mathcal F)\cdot\alpha<0.
\]
The bracket obstruction from $\bigwedge^2\mathcal F$ to
$T_X/\mathcal F$ vanishes by the slope inequality.  Thus $\mathcal F$
is a foliation, and the cited algebraicity theorem realizes it as the
relative tangent foliation of a rational map with connected general
fibers.  Prepare an equidimensional model of that map, then take a
smooth resolution with maps $\mu:W\to X$ and $f:W\to Y$.
Put $L=-K_X$ and retain its smooth semipositive metric.

The relative adjoint bundle $K_{W/Y}+\mu^*L$ has a nonzero section on
a general fiber: there it is represented by the effective Jacobian
$K_{W/X}$.  On the smooth bundle and base-change locus, evaluation
with the $L^2$ direct-image metric supplies its Bergman weight
\[
 \Phi(x)=\log\sup_{\|u\|_{L^2}\le1}|u(x)|^2
\]
in a local relative-adjoint frame.  Semipositivity follows from the
dual-norm form of direct-image positivity
\cite{Berndtsson2009,Berndtsson2011}; zeros of the evaluation map give
allowed singularities.  On the isomorphism locus of $\mu$, the
submersive differential identifies this line with
$\det(T_X/\mathcal F)$.  The only extra codimension-one issue is
ramification.

A prime of $X$ vertical for the map dominates a base prime on the
normal equidimensional model.  Near its general point, and hence on
its strict transform in $W$, choose coordinates
\[
 f=(z_1^e,z_2,\ldots,z_{\dim Y}).
\]
The differential contributes $z_1^{e-1}$.  Therefore
$\det(T_X/\mathcal F)$ corresponds to the relative adjoint line
twisted down by $(e-1)\{z_1=0\}$.
A relative form written in total-space/base canonical frames acquires
the factor $(ez_1^{e-1})^{-1}$ when restricted to a fiber.
Integration on a fixed local fiber chart and the holomorphic submean
estimate consequently bound every unit-norm section coefficient by
$C|z_1|^{e-1}$.  It follows that
\[
 \Phi\le(e-1)\log|z_1|^2+O(1).
\]
After the indicated twist, the induced weight is locally bounded
above at the general point of every such prime.  At horizontal primes
the ordinary submean estimate gives the same upper boundedness without
a ramification term.  These weights extend plurisubharmonically across
codimension one, and then across the remaining codimension-two locus
on $X$.  They define a semipositive singular metric on
$\det(T_X/\mathcal F)$, making that line pseudoeffective.  This
contradicts its negative degree against $\alpha$.

By movable-cone duality, the negative determinant of every cotangent
tensor subsheaf is therefore pseudoeffective.  The argument retains the
factor $e-1$ because omitting it would construct a metric on the wrong
line at ramified divisors.
